%% ===========================================================================
%%  Bab 3 Profil Organisasi dan Aktivitas Magang
%% ===========================================================================
\chapter{Profil Organisasi dan Aktivitas Magang}

\section{Gambaran Umum Organisasi}
% TODO: profil \InstansiMagang: sejarah singkat, bidang kerja, struktur
% organisasi, dan penempatan unit tempat magang.

\section{Aktivitas Magang}
% TODO: uraikan aktivitas selama magang, boleh dibagi per pekan, per proyek,
% atau per tahap pekerjaan. Sebutkan periode \PeriodeMagang.

\section{Metode Pelaksanaan Magang}
% TODO: cara pelaksanaan kerja, misalnya daring/luring, bimbingan berkala,
% rapat tim, dan alur pelaporan.

\section{Bentuk-bentuk Dukungan}
% TODO: satu alinea pengantar, lalu rincian pada tiga subbagian di bawah.

\subsection{Dukungan dari Mahasiswa}
% TODO: peran, inisiatif, dan kesiapan mahasiswa.

\subsection{Dukungan dari Institusi Tujuan Magang}
% TODO: fasilitas, mentor, data, dan akses yang diberikan instansi.

\subsection{Dukungan dari Program Studi}
% TODO: dukungan administratif, pembimbing, dan kebijakan kampus.

\section{Kendala-kendala}
% TODO: satu alinea pengantar, lalu rincian pada tiga subbagian di bawah.

\subsection{Kendala dari Mahasiswa}
% TODO: kendala teknis dan nonteknis dari sisi mahasiswa.

\subsection{Kendala dari Institusi Tujuan Magang}
% TODO: kendala dari sisi instansi, misalnya akses data atau perizinan.

\subsection{Kendala dari Program Studi}
% TODO: kendala dari sisi kampus, misalnya administrasi atau jadwal.
